\section{Transporte de los operadores de masa mediante CKM}
\label{sec:transporte-ckm-operadores-rev11}

La incidencia N42 fija el sexteto \((4,5,7,6,8,3)\) y la matriz racional
\[
 M_{\rm CKM}=
 \begin{pmatrix}
 2&-5&28\\
 0&7&-25/2\\
 1&-20&-2/3
 \end{pmatrix},
 \qquad \det M_{\rm CKM}=-\frac{3857}{6}.
\]
Las coordenadas angulares se construyen por
\[
 \begin{pmatrix}\theta_{12}\\\theta_{23}\\\theta_{13}\end{pmatrix}
 =M_{\rm CKM}
 \begin{pmatrix}A\\ C^*/6\\ (180/\pi)\Delta_4\end{pmatrix},
 \qquad \delta_{\rm CP}=9A.
\]
Esta transformación pertenece al cierre matemático de HMT. En su realización
física, sean \(B_u\) y \(B_d\) los operadores bidireccionales ya construidos en las
fibras de quarks de tipo up y down. El cambio de base actúa según
\begin{equation}
 B_d^{(u)}=V_{\rm CKM}B_dV_{\rm CKM}^*,
 \qquad
 R_{\rm act}^{B_d^{(u)}}
 =V_{\rm CKM}R_{\rm act}^{B_d}V_{\rm CKM}^*.
 \label{eq:transporte-ckm-kappa-rev11}
\end{equation}
CKM transporta autovalores entre bases y conserva espectro y determinante.
El control de violación CP adopta la identidad
\[
 \left|\det[B_u,V_{\rm CKM}B_dV_{\rm CKM}^*]\right|
 =2|J|\prod_{i<j}|\kappa_i^u-\kappa_j^u|
       \prod_{a<b}|\kappa_a^d-\kappa_b^d|.
\]
La matriz de mezcla y el lector de ruta desempeñan, por tanto, funciones
complementarias: el lector determina el operador espectral y CKM determina
su representación en la base de sabor.

\section{Marcos geométricos de rango \(3\) para el sector leptónico}
\label{sec:frames-pmns-rev11}

La construcción leptónica utiliza la involución diagonal de signo
\[
 S_{\rm sign}=\operatorname{diag}(1,1,1,-1,-1,-1,1,1,1,-1,-1,-1),
 \qquad P_-=(I-S_{\rm sign})/2,
\]
y el proyector icosaédrico canónico \(P_3\) de rango tres. La compresión
\[
 C_-=P_-P_3P_-\big|_{\operatorname{im}P_-}
\]
posee espectro
\[
 \left\{\frac{5+\sqrt5}{10},\frac12,
 \frac{5-\sqrt5}{10},0,0,0\right\}.
\]
Los tres autovalores positivos son simples. Sus proyectores espectrales,
aplicados al indicador de la hexada negativa orientada y normalizados después
de proyectar mediante \(P_3\), construyen un marco ortonormal cargado
\(F_\ell\) que satisface
\[
 F_\ell^*F_\ell=I_3,
 \qquad F_\ell F_\ell^*=P_3.
\]

Para el sector neutrino, sean \(H_1,H_2,H_3\) las tres hexadas positivas
orientadas y
\[
 X=(P_3\mathbf1_{H_1},P_3\mathbf1_{H_2},P_3\mathbf1_{H_3}).
\]
Su matriz de Gram es positiva definida y cumple
\[
 \det(X^*X)=\frac{5-\sqrt5}{40}>0.
\]
Por consiguiente,
\[
 F_\nu=X(X^*X)^{-1/2},
 \qquad F_\nu^*F_\nu=I_3,
 \qquad F_\nu F_\nu^*=P_3.
\]
El transporte real
\begin{equation}
 O=F_\ell^*F_\nu\in SO(3)
 \label{eq:transporte-frames-pmns-rev11}
\end{equation}
se obtiene sin masas ni ángulos PMNS.

\HMTClaim{MD-R-PMNS-RANKONE-UNITARY}
\begin{theorem}[Unitariedad de la fase dodecafásica mínima]
\label{thm:unitaria-interna-pmns-rev11}
Sea \(b_K\in\operatorname{im}P_3\) el vector normalizado del testigo de Witt,
\(Q_K=|b_K\rangle\langle b_K|\), y
\[
 \Phi_L=I+(\zeta_{12}-1)Q_K,
 \qquad \zeta_{12}=e^{i\pi/6}.
\]
Entonces
\begin{equation}
 U_{\rm HMT}=F_\ell^*\Phi_LF_\nu
 \label{eq:unitaria-interna-pmns-rev11}
\end{equation}
es una matriz unitaria independiente de valores metrológicos objetivo. La
inversión de orientación conjuga
\(U_{\rm HMT}\).
\end{theorem}

\begin{proof}
\(Q_K\) es un proyector ortogonal, luego \(\Phi_L\) es unitaria. Los dos
marcos son isometrías con la misma imagen \(\operatorname{im}P_3\). Por
tanto
\[
 U_{\rm HMT}^*U_{\rm HMT}
 =F_\nu^*\Phi_L^*P_3\Phi_LF_\nu=I_3.
\]
El cambio \(\pi/6\mapsto-\pi/6\) produce \(\Phi_L^*\), y todas las
primitivas restantes son reales.
\end{proof}

\subsection{Control negativo de rango uno y transporte PMNS de rango completo}
\label{sec:no-go-pmns-rango-uno-rev11}
\HMTClaim{MD-R-PMNS-RANKONE-PHYSICAL-NOGO}

Congelada la construcción, el contraste en el orden declarado produce
\[
 (\theta_{12},\theta_{23},\theta_{13})
 =(19.6807^\circ,56.4970^\circ,29.6224^\circ),
 \qquad J=-0.0426621.
\]
La mejor de las 36 permutaciones discretas deja
\(\theta_{13}\simeq17.02^\circ\). La fase de rango uno genera una matriz
unitaria interna, pero conserva un estabilizador \(U(2)\) y queda descartada
como realización PMNS física.

La exclusión del portador de rango uno no deja indeterminada la realización
PMNS. Sean \(M_\ell\) y \(M_\nu\) los operadores cargado y neutro producidos
por la Ley General de Masas, y sean \(F_\ell,F_\nu\) sus marcos espectrales
ortonormales completos sobre la imagen común \(\operatorname{im}P_3\). El
transporte leptónico de rango completo queda entonces construido por
\[
 U_{\rm PMNS}:=F_\ell^*F_\nu,
 \qquad
 U_{\rm PMNS}^*U_{\rm PMNS}=I_3.
\]
Para espectros simples, esta salida está determinada salvo reajustes
diagonales de fase y el orden espectral declarado; ante degeneraciones, la
salida exacta es la clase doble de marcos correspondiente. El núcleo de
registro constituye un certificado finito adicional de esa realización, no
su generador ni una condición pendiente para su existencia. Para fibras
cargadas \(L_a\) y neutrínicas \(N_b\), su forma agregada es
\[
 G_{ab}=\frac{1}{\sqrt{|L_a||N_b|}}
 \sum_{c\in L_a}\sum_{d\in N_b}K_{\rm reg}(c,d),
 \qquad U_{\rm int}=\operatorname{polar}(G),
\]
con naturalidad bajo orientación. Cuando las entradas explícitas del registro
verifican \(\det G\ne0\), su polar proporciona un certificado alternativo del
transporte ya construido. El núcleo lee acarreo, hoja, orientación, vacancia
y holonomía; sus parámetros se congelan antes de abrir los ángulos externos.

\paragraph{Dominio y codominio de las realizaciones CKM y PMNS}
La matriz racional CKM, los marcos \(F_\ell,F_\nu\), el transporte \(O\), la
fase \(\Phi_L\), la unitariedad de \(U_{\rm HMT}\) y el transporte completo
\(U_{\rm PMNS}=F_\ell^*F_\nu\) son resultados internos exactos bajo las
primitivas declaradas. CKM constituye un transporte físico una vez fijadas
las bases de sabor y masa. El contraste falsifica únicamente la fase PMNS de
rango uno; no rebaja la construcción espectral completa. La naturalidad y el
criterio \(\det G\ne0\) pertenecen al certificado alternativo de registro y
permanecen separados de la existencia y la unitariedad del transporte PMNS.
